\section{Conclusion}
This thesis took a notebook comparison of five classifiers on two Twitter corpora and re-built it as a 56-cell benchmark under one leakage-free protocol, with three corpora, five representations and eight model families, delivered as a system that tracks its runs, documents and serves its models and monitors them in production. The protocol changed the conclusions more than the new models did. On short imbalanced text all families rank documents almost equally well, and the differences that fixed-threshold comparisons report are differences in calibration and threshold handling; tuning the threshold out-of-fold is worth up to 0.14 F1 and the least calibrated models transfer it worst. Sparse $n$-gram features remain the representation of choice for classical learners, learned document embeddings are the weakest choice everywhere, and on balanced long documents linear models on TF-IDF match networks trained from scratch at a hundredth of the cost. The original study's ranking survives this scrutiny; its numbers do not.

Every number in this document is regenerated from saved predictions by the public code, and the selected models can be served and monitored from the same repository. That, more than any individual result, is what the 2022 study lacked.

\section{Future work}
The results and the state of the field in 2026 (Section~\ref{sec:bg2026}) define a research programme. The items are ordered by how directly they follow from the evidence and by cost.

\subsection{Cross-validation bagging and calibrated decisions}
The one clear defect the protocol exposed is threshold transfer for over-confident models. The direct remedy is to serve the average of the $k$ fold models rather than a model refitted on all the data: the served probabilities are then on the same scale as the out-of-fold probabilities on which the threshold was chosen, at the cost of $k$ times the inference. A second remedy is post-hoc calibration of the refitted model on the out-of-fold predictions (temperature scaling for the networks, isotonic regression for the trees). Both are cheap experiments on the saved probabilities and should be the first extension. Beyond a single threshold, conformal prediction~\cite{angelopoulos2023} turns any classifier's scores into prediction sets with a guaranteed coverage rate on exchangeable data; for a moderation setting, where the cost of a missed positive differs from that of a false alarm, sets with guaranteed recall are a more honest interface than a point decision at a tuned threshold.

\subsection{Pretrained encoders as the upper reference}
Fine-tuned transformer encoders~\cite{devlin2019,liu2019} should be added to the grid as a reference family, with the same protocol: the threshold chosen out-of-fold, the test split scored once. Parameter-efficient fine-tuning~\cite{hu2022} makes this feasible on a single GPU within the hour, and the calibration analysis of Chapter~\ref{ch:results} predicts that such models will need temperature scaling before their thresholds transfer. The value of the experiment is not the higher number, which is known, but the measurement of how much of the gap to the linear models is closed on the hate corpus, where label noise may cap every model, against IMDB, where the literature suggests seven points of accuracy.

\subsection{Large language models as annotators and as classifiers}
The vaccine corpus is the weakest link of the study because its labels are a lexicon. Two experiments follow. First, annotate a stratified sample of one to two thousand tweets with human judges and measure the agreement of TextBlob and of the trained models with the humans, which would tell whether the 0.81 macro F1 of the CNN means anything. Second, following Gilardi et al.~\cite{gilardi2023} and the survey of Zhang et al.~\cite{zhang2023}, obtain labels from an instruction-following large language model with a fixed prompt and measure its agreement with the human sample; if it is high, relabel the whole corpus and rerun the grid. The same models can be evaluated as zero-shot classifiers within the protocol, with their verbalised confidence treated as a score to be calibrated like any other. The comparison of cost per label between crowd workers, a large model and a fine-tuned small model is itself a result worth reporting.

\subsection{Fairness auditing of hate-speech classifiers}
Sap et al.~\cite{sap2019} showed that hate-speech classifiers trained on Twitter corpora flag tweets in African-American English at higher rates. A classifier selected by F1 alone can be selected into this bias. The benchmark should add a fairness slice: a dialect or demographic proxy for the hate corpus, per-group false-positive rates at the tuned threshold, and a selection rule that constrains the gap. The model card already has a place for such results; the protocol does not yet produce them.

\subsection{Drift-triggered retraining and continual evaluation}
The monitor detects drift in input length, vocabulary and predicted-class distribution; it does not yet act on it. The next step is a retraining pipeline triggered by an alert: collect the drifted traffic, obtain labels for a sample (by the annotation route above), rerun the benchmark on the union of old and new data, compare the new selected model with the deployed one on the new sample under the same protocol, and promote it only if it wins. Each promotion should leave a model card and a tracked run behind. Over time this yields a longitudinal record of how the task itself changes, which no static benchmark can give.

\subsection{Broader coverage}
The grid should grow along four axes. Corpora: a human-annotated hate-speech corpus with published guidelines~\cite{davidson2017}, a multilingual corpus to test whether the representation findings hold outside English, and a time-split corpus to measure temporal drift directly. Representations: pretrained static embeddings as a control between Word2Vec trained on the fold and contextual encoders. Models: a character-level model for the misspelled and hashtag-laden hate corpus, and NB-SVM~\cite{wang2012} as the strongest known linear baseline on IMDB. Analysis: per-cell hyper-parameter search inside a nested loop, to measure how much of each family's gap to the best is due to untuned settings rather than to the family.

\subsection{Cost and energy accounting}
Training time was reported because it was free to measure. Energy and monetary cost per training run and per thousand predictions should be recorded with the same rigour, including the pretrained and large-model references; the efficiency frontier of Chapter~\ref{ch:results} is the figure most decision makers actually need, and in 2026 its horizontal axis is increasingly measured in currency rather than seconds.

\section{Closing remark}
The most useful outcome of this work is not a ranking of classifiers but a demonstration that the ranking depends on the evaluation, and a piece of software that makes the right evaluation cheaper than the wrong one. The programme above applies that software to the questions that the next five years of sentiment analysis will actually turn on: what large models make possible, what they cost, and how to tell, without fooling oneself, whether they are better.
